\subsection{范数收敛}
\label{sec:ch1-norm-convergence}

测度论和 $L^p$ 空间的发展为收敛问题带来了新的研究途径. 现在可以提出:
\begin{enumerate}[label=(\arabic*)]
\item 对 $f\in L^p(\mathbb T)$, 是否有 $\lim_{N\to\infty}\|S_Nf-f\|_p=0$?
\item 对 $f\in L^p(\mathbb T)$, 是否有 $\lim_{N\to\infty}S_Nf(x)=f(x)$ 几乎处处成立?
\end{enumerate}
第一个问题可用下述引理重新表述.

\begin{Lemma}
\label{lem:ch1-lp-partial-sum-convergence}
当 $1\le p<\infty$ 时, $S_Nf$ 在 $L^p$ 范数下收敛到 $f$, 当且仅当存在与 $N$ 无关的 $C_p$, 使得
\begin{equation}
\label{eq:ch1-partial-sum-lp-bound}
 \|S_Nf\|_p\le C_p\|f\|_p.
\end{equation}
\end{Lemma}

\begin{Proof}
\eqref{eq:ch1-partial-sum-lp-bound} 的必要性由一致有界原理得到.

为证明其充分性, 先注意, 若 $g$ 是三角多项式, 则当 $N\ge\deg g$ 时 $S_Ng=g$. 因此, 由三角多项式在 $L^p$ 中稠密(见推论~\ref{cor:ch1-trigonometric-density-uniqueness}), 对 $f\in L^p$, 可以找到三角多项式 $g$, 使得 $\|f-g\|_p<\varepsilon$. 所以, 当 $N$ 充分大时,
\[
 \|S_Nf-f\|_p
 \le\|S_N(f-g)\|_p+\|S_Ng-g\|_p+\|f-g\|_p
 \le(C_p+1)\varepsilon.
\]
\end{Proof}

若 $1<p<\infty$, 则不等式 \eqref{eq:ch1-partial-sum-lp-bound} 成立, 我们将在 \MBUnresolvedReference{chapter}{第 3 章} 证明这一点. 当 $p=1$ 时, $S_N$ 的 $L^1$ 算子范数仍为 $L_N$, 所以由引理~\ref{lem:ch1-lebesgue-number-asymptotic}, 第一个问题的答案是否定的.

当 $p=2$ 时, 由 \eqref{eq:ch1-exponential-orthogonality}, 函数系 $\{e^{2\pi ikx}\}$ 是正交归一系; 又由三角多项式在 $L^2$ 中稠密, 它是完备的(即为正交归一基). 因而可应用 Hilbert 空间理论得到下述结果.

\hypertarget{chapter-01-page-021}{}
\begin{Theorem}
\label{thm:ch1-fourier-l2-isometry-series}
映射 $f\mapsto\{\widehat f(k)\}$ 是从 $L^2$ 到 $\ell^2$ 的等距映射, 即
\[
 \|f\|_2^2=\sum_{k=-\infty}^{\infty}|\widehat f(k)|^2.
\]
\end{Theorem}

$L^2$ 中的范数收敛立即由此得到.

第二个问题困难得多. A. Kolmogorov 于 1926 年给出了一个可积函数, 其 Fourier 级数在每一点都发散, 所以当 $p=1$ 时答案是否定的. 若 $f\in L^p$, $1<p<\infty$, 则 $f$ 的 Fourier 级数几乎处处收敛. 这一点由 L. Carleson 于 1965 年对 $p=2$ 证明, 由 R. Hunt 于 1967 年对 $p>1$ 证明. 在 Carleson 的结果出现之前, 即使对连续函数 $f$, 这一问题的答案也是未知的.

